\section{系统运行与维护}

\subsection{部署}

系统以 Docker 镜像发布。推送到主分支后，持续集成先在两种数据库上运行测试，再构建镜像并推送到镜像仓库。服务器上的 Docker Compose 启动三个容器：Web 服务、后台任务和机器人。三者共享同一数据库，用量库以只读方式挂载。升级步骤固定为“备份 → 拉取新镜像 → 用新镜像执行 \code{alembic upgrade head} → 重启服务”，迁移成功后才会切换到新版本的服务。

数据库的选择通过环境变量决定：
\begin{itemize}
  \item 未设置 \code{BILLING\_DATABASE\_URL} 时，使用数据目录下的 SQLite 文件，这是单机部署的默认方式；
  \item 设置为 \code{postgresql+psycopg://…} 时，连接 PostgreSQL。此时服务应以 \code{billing\_app} 角色连接，迁移以属主角色执行。
\end{itemize}
服务每次启动都会检查触发器和视图是否齐全，只补建缺失的对象。PostgreSQL 上用会话级咨询锁（advisory lock）避免三个进程同时执行 DDL；SQLite 上用 \code{IF NOT EXISTS} 保证并发启动的安全。

\subsection{备份与恢复}

\begin{itemize}
  \item \textbf{备份}：执行 \code{cpa-billing backup <目标文件>}。SQLite 使用在线备份接口，PostgreSQL 调用 \code{pg\_dump --format=custom}。两种方式都能在服务运行时得到一致的快照，可以配合 cron 每日执行，并保留最近 30 份。
  \item \textbf{SQLite 恢复}：停止三个服务，用备份文件替换数据库文件（同时删除旧的 \code{-wal} 和 \code{-shm} 文件），然后重新启动。
  \item \textbf{PostgreSQL 恢复}：在新的空库上执行 \code{pg\_restore --no-owner -d billing <备份文件>}，再执行一次 \code{roles.sql} 恢复授权。
  \item \textbf{升级前回退点}：每次升级前都要备份。迁移脚本提供 \code{downgrade}，例如 0017 的回退会删除触发器、视图和部分唯一索引，但不会删除任何业务数据。
\end{itemize}

\subsection{日常运维与监控}

\begin{itemize}
  \item \textbf{同步状态}：管理控制台显示同步检查点的积压条数、最近一次成功时间和最近一次错误；后台任务每轮都会在日志中记录导入和计价条数。
  \item \textbf{死信处理}：用量库中无法解析的记录进入死信表，后台任务每轮会自动重试；管理控制台可以查看仍未解决的死信。
  \item \textbf{对账}：\code{cpa-billing reconcile} 比较用量库、原始记录和计价记录三者的条数，统计未归属、未分配资源池和未计价的事件，并把结果写入对账记录表。账期关账前要求未计价事件为 0，否则必须填写数据质量豁免说明。
  \item \textbf{审计}：所有管理操作都记录操作者、对象以及修改前后的值，审计角色可以独立查询。
  \item \textbf{数据增长}：调用记录每月约增长数十万行。当单表达到数千万行时，可以按发生时间对调用记录和计价记录做范围分区，并把已关账账期的分区设为只读。现有索引都以时间为主要或次要列，可以平滑过渡到分区表。
\end{itemize}

\subsection{总结}

本课程设计在一个真实使用中的共享计费系统上，按软件工程流程完成了从问题定义、需求分析到数据库设计、实现、测试和维护的全过程。数据库设计方面的主要工作和体会如下：

\begin{enumerate}
  \item \textbf{用版本化和快照保证账单可复现。}价格、归属、规则都会变化，这是需求最本质的复杂性；价格版本、归属期和账期快照的建模决定了“每期账单可查询、可核对”能否成立。
  \item \textbf{把不变量交给数据库强制执行。}约束、触发器和角色权限在数据库层再落实一遍应用层规则，才能防住手工 SQL 和未来代码的疏漏；设计时先通读写路径，再用全部测试回归，避免误伤正常流程。
  \item \textbf{3NF 是起点而不是教条。}保留的快照属性是历史事实，冗余列有明确的性能收益，二者的一致性由触发器保证。
  \item \textbf{性能结论以测量为准。}百万级测试既验证了索引，也表明账期全量聚合等场景应靠减少计算次数和增量处理来优化。
  \item \textbf{可移植性需要测试保证。}移植暴露的 10 类方言差异中，整数除法的舍入差异不会报错、只会让统计错位，只有双库运行同一套测试才能发现。
\end{enumerate}

后续可以改进的方向包括：按时间对大表分区；为成员提供账单异议与确认流程；把对账任务纳入定时调度并在异常时自动告警。
